\section{Behavioural evaluation of LLM agents}\label{sec:agentbench}

The agent-native design described in Section~\ref{sec:agent} is the
most contestable claim in this paper: a reviewer entitled to
scepticism may reasonably ask whether a machine-readable schema
\emph{actually} improves the behaviour of an LLM agent on a real
empirical task, or whether the \emph{statsmodels-plus-prompting}
baseline performs equivalently well. We treat that question as
falsifiable. This section reports the
\emph{mechanical} and \emph{procedural} evidence we can ship with
the present draft, sketches the \emph{behavioural} evidence we have
prepared for production, and defers the production RCT itself to
a forthcoming companion benchmark paper to preserve its
pre-registration value. We refer to \citet{patil2023gorilla} and
\citet{liang2023helm} as methodological precedents for tool-use and
holistic-LLM benchmarks respectively.

\subsection{Mechanical evidence}\label{sec:agentbench-mechanical}

The function-schema layer is statically inspectable. We report in
Table~\ref{tab:schema-coverage} that 980 of 980 \statspai{} public
functions return a non-empty OpenAI-style JSON Schema; that 214
functions have hand-written registry specifications; and that 78
functions currently expose richer agent metadata cards with
assumptions, pre-conditions, failure modes, and recovery hints. The
closest comparable \proglang{Python} statistics packages
(\pkg{statsmodels}, \pkg{linearmodels}, \pkg{scikit-learn}) expose
zero package-wide typed tool schemas.

\begin{table}[t]
\centering\small
\caption{Mechanical schema coverage at \statspai{} 1.14.0 against
representative comparators. Counts refer to public functions or
estimator classes.}
\label{tab:schema-coverage}
\begin{tabular}{lcccc}
\toprule
& \statspai{} & \pkg{statsmodels} & \pkg{linearmodels} & \pkg{scikit-learn} \\
\midrule
Public surface (functions/classes)              & 980 & $\sim$2{,}000 & $\sim$150 & $\sim$170 estimators \\
Emits OpenAI-style tool schema                  & 980 & 0   & 0   & 0 \\
Hand-written registry specification             & 214 & 0   & 0   & 0 \\
Curated agent metadata card                     & 78  & 0   & 0   & 0 \\
Known limitations surfaced before execution     & 10  & 0   & 0   & 0 \\
\bottomrule
\end{tabular}
\end{table}

\subsection{Procedural evidence}\label{sec:agentbench-procedural}

The agent-native surface is exposed through an MCP server
(Section~\ref{sec:mcp}) that is conformant to the open Model Context
Protocol specification. Conformance is asserted by the local
\code{tests/test\_mcp\_protocol.py}, \code{tests/test\_agent\_schema.py},
and \code{tests/test\_registry.py} suites: they check that registered
functions emit tool schemas, that agent-card filters behave as
documented, and that MCP tool discovery sees the same registry-backed
surface as \code{sp.function\_schema()}. This guarantees that the
schema layer is, at a minimum, correctly wired up before any
behavioural claim is made.

For this draft we also execute one deterministic end-to-end trace
through the live MCP server. The script
\code{replication/scripts/ex07\_agent\_trace.py} writes a temporary
\texttt{mpdta.csv}, calls \code{tools/list}, fits
\code{callaway\_santanna} with \code{as\_handle=true}, sends the
returned handle to \code{audit\_result}, and then sends the same
handle to \code{honest\_did\_from\_result}. Table~\ref{tab:agent-trace}
summarises the trace; Appendix~\ref{app:transcript} prints the
normalised transcript. This is not a substitute for the production
RCT, but it verifies the agent-specific claim that the estimate
\(\to\) audit \(\to\) sensitivity chain is executable without
ferrying arrays or hand-copying citations through the model context.

\begin{table}[t]
\centering\small
\caption{Deterministic MCP trace generated by
\code{replication/scripts/ex07\_agent\_trace.py}. Ephemeral result
handles are normalised in the appendix transcript.}
\label{tab:agent-trace}
\input{../replication/tables/ex07_agent_trace.tex}
\end{table}

\subsection{Behavioural evidence (deferred to companion paper)}\label{sec:agentbench-behavioural}

Mechanical and procedural evidence do not by themselves answer
whether an LLM agent invoked through this interface produces
\emph{better empirical analyses} than one invoked against
\pkg{statsmodels} or against \proglang{R} packages exposed through a
Jupyter MCP shim. The benchmark that closes this gap --
\textsc{CausalAgentBench} -- is an RCT-style study that holds the
agent's language model fixed and varies only the toolset.

The protocol consists of fifty causal-inference research prompts
distributed across three difficulty levels (twenty L1 method-named
prompts, twenty L2 indirect prompts, ten L3 workflow prompts), six
experimental cells in a $3 \times 2$ factorial design (\statspai{}
with MCP, a Pythonic statsmodels/linearmodels/DoubleML/grf-python
stack, and \proglang{R}-via-MCP, each crossed with a frozen
Anthropic Claude release and a frozen OpenAI GPT release), and a
total of 900 trials at three repetitions per (prompt, cell). Eight
metrics span task success, method correctness, code-execution
success, token efficiency, hallucination rate, diagnostic
completeness, reproducibility, and time-to-result. Five hypotheses
will be pre-registered on OSF; failed trials are classified into
five qualitative failure modes and reported with redacted
transcripts. The cluster bootstrap with prompt as the clustering
unit, $\alpha = 0.05$ two-sided with Bonferroni correction across
the five hypotheses, is the planned statistical test.

The complete protocol, gold answers, grading rubric, and a
deterministic mock-LLM harness are shipped at
\code{tests/agent\_bench/} in the source tree; an OSF
pre-registration draft frozen at
\code{tests/agent\_bench/prompts/\_protocol.md} lists the file
hashes to deposit before the first production trial. A 900-trial
mock-LLM dry run completes in $<1$~s on the benchmarking machine,
validates the harness end-to-end, and produces the same per-cell
table format that the production run will report. Flipping
\code{runner.py} from \code{--mock} to \code{--api both} is the
single change that swaps the mock LLM for the real ones; the
scoring, aggregation, and statistical-test pipeline downstream of
that flip is identical. The production run is gated on two
external steps that we deliberately do not auto-execute: the OSF
pre-registration deposit and an API budget approval. The full RCT
and its statistical analysis are the subject of a forthcoming
companion benchmark paper; we defer reporting them here both to
preserve the pre-registration's scientific value and to keep the
present manuscript focused on the unification, parity, and
schema-layer claims that constitute the core JSS contribution.
